\section{Tamamlayıcı vaka: örneklem dosyası çerçeve değildir}
\label{sec:v2-nhanes}

\textbf{Soru ve hedef.} NHANES 2017--2018 araştırmasında ABD'nin 50
eyaleti ve Columbia Bölgesi'nde yaşayan, kurumlarda yaşamayan sivil
nüfus hedeflenir. Bu vaka için sorumuz daha dardır: bu nüfusun 20 yaş
ve üzerindeki kesiminde, 60 yaş ve üzerindekilerin oranı nedir?
Dönem ve yaş sınırları analiz tanımının parçasıdır; hedef bugünkü
nüfus veya Türkiye nüfusu değildir. Bu bir müdahale etkisi sorusu da
değildir. Veri kaynağı demografi dosyası \texttt{DEMO\_J}'dir
\citep{nchs2020demoj,nchssampledesign}.

\textbf{Çerçeveden kayda.} Seçim, coğrafi birincil örnekleme
birimlerinden segmentlere, hanelere ve kişilere ilerleyen çok aşamalı
bir düzene dayanır. Hane listeleri ve tarama işlemleri, kişilerin
seçilebilmesini sağlar. Kamuya açık yanıtlayıcı dosyası bu çerçevelerin
yerine geçmez: seçilmeyen veya yanıt vermeyen bütün birimleri içermez
\citep{nchssampledesign}. B06'daki yapay çerçeveden bilgisayarla seçim
yapmak, burada uygulanmış saha tasarımını yeniden gerçekleştirmek değildir.

\textbf{Analiz tanımı.} Gözlem birimi kişidir. Yaş için
\texttt{RIDAGEYR}, görüşme ağırlığı için \texttt{WTINT2YR}; varyans
hesabı için \texttt{SDMVSTRA} ve \texttt{SDMVPSU} kullanılacaktır.
Son iki değişken gizliliği koruyan varyans birimleridir; gerçek coğrafi
yer adları veya doğrudan saha çerçevesi olarak okunmaz. Yaştaki 80
kodu 80 ve üzerini birleştirir; seçtiğimiz 60 yaş eşiğini belirsiz
hale getirmez \citep{nchs2020demoj}.

\begin{definitionbox}
\textbf{İki farklı özet.} $d_i=\Ind{\text{yaş}_i\geq20}$ ve
$y_i=\Ind{\text{yaş}_i\geq60}$ olsun. $w_i$ görüşme ağırlığıdır.
\[
\widehat p_w=\frac{\sum_i w_i d_i y_i}{\sum_i w_i d_i},
\qquad
\widehat p_0=\frac{\sum_i d_i y_i}{\sum_i d_i}.
\]
İlk ifade hedef alt nüfusun oranını tahmin etmek için kullanılacak
ağırlıklı orandır. İkincisi dosyadaki yetişkin yanıtlayıcıların
bileşimini betimler. İki sonuç karşılaştırılabilir; fakat tercih edilen
sayısal sonucu verdiği için ağırlık seçilmez veya kaldırılmaz.
\end{definitionbox}

\textbf{Ağırlık neyi düzeltir?} Nihai ağırlık, temel seçim ağırlığının
yanı sıra yanıtsızlık ve nüfus toplamlarına uyarlama işlemlerini de
içerir. Bu nedenle $1/w_i$ değerini özgün seçilme olasılığı diye
raporlamayız. Yalnız görüşme değişkenlerini kullanan bu soru için
muayene ağırlığına geçmeyiz. Ağırlıklandırma, çerçeve kapsamı veya
yanıtsızlıktan doğan her yanlılığın giderildiğini garanti etmez
\citep{nchsweighting}.

\textbf{Belirsizlik nasıl hesaplanacak?} Tabaka, küme ve ağırlıkları
birlikte kullanan Taylor doğrusallaştırmasıyla tasarıma uygun standart
hata hesaplanmalıdır. Yetişkinler bir \emph{alt evren} olarak tanımlanır;
tasarım kurulmadan çocuk kayıtları silinmez. Yalnız ağırlıklı oranı
bulmak veya $\sqrt{\widehat p(1-\widehat p)/n}$ kullanmak bu hesabın
yerine geçmez \citep{nchsvariance}. Tahmin ile güven aralığı, farklı
tamamlanma kontrolleridir.

\begin{researchbox}
\textbf{Vakanın mevcut kapsamı.} Kaynak belgelerine dayalı tasarım
okumasını, görevlerin ardından gelen sayısal uygulama tamamlar.
Dosya kimliği, yaş kodları, pozitif ağırlıklar, tabaka--PSU yapısı ve
seçili alanlarda eksiksizlik kontrol edilmiştir. Diğer değişkenlerdeki
yapısal eksiklikler yüzünden bütün sütunlarda topluca tam kayıt
seçilmemiştir. Beklenmedik girdide kod otomatik dışlama yerine durur.
Bu dar yaş oranı uygulaması NHANES'in bütün analizleri için kabul
veya yayın onayı değildir.
\end{researchbox}

\subsection{Karar görevleri ve gerekçeli kontrol}
Önce aşağıdaki beş görevi yanıtlayın; ardından kontrol yanıtlarıyla
karşılaştırın. Sayısal veri dosyası olmadan da tasarım kararları
değerlendirilebilir.

\begin{enumerate}
\item Hedef nüfusu, alt evreni ve dönemi yazın. Kurumlarda yaşayanları
aynı hedefe dahil etmek neden yalnız ağırlık değişikliği değildir?
\item Bir öğrenci ``\texttt{DEMO\_J}'den rastgele 200 satır seçersem
özgün örnekleme çerçevesini elde ederim'' diyor. Yanılgıyı açıklayın.
\item İki oran formülünün paydalarını yorumlayın. Nihai ağırlığın tersi
neden tek başına saha seçiminin olasılığını vermez?
\item ``Çocukları siler, ağırlıklı oranı bulur, bağımsız gözlem
formülüyle standart hata hesaplarım'' planını düzeltin.
\item Sonuç raporu için iki kapsam cümlesi yazın: biri kime ve hangi
döneme ilişkin tahmin yapılacağını, diğeri hangi iddianın
kurulamayacağını belirtsin. Henüz hesaplanmamış değer yazmayın.
\end{enumerate}

\textbf{Kontrol yanıtları.}
\begin{enumerate}
\item Araştırmanın hedefindeki 20 yaş ve üzeri kişiler alt evrendir;
dönem 2017--2018'dir. Kurum nüfusunu eklemek, hedefi ve gerekli kapsamı
değiştirir; mevcut yanıtlayıcılara yeni katsayı vermek yeterli değildir.
\item Dosya, seçim ve yanıt süreçlerinden sonraki kayıttır. Ondan
alt örnek seçmek, seçilmeyen haneleri veya yanıt vermeyen kişileri
geri getirmez; ayrıca yeni bir seçim aşaması yaratır.
\item Ağırlıklı payda yetişkin alt nüfusun büyüklüğüne ilişkin tahmin;
ağırlıksız payda yetişkin kayıt sayısıdır. Nihai ağırlıktaki ek
düzeltmeler nedeniyle temel $1/\pi_i$ ile son ağırlık ayrılmalıdır.
\item Önce tam tasarım bilgisi korunur, sonra yetişkin alt evreni
tanımlanır. Varyans hesabı tabaka ve PSU katkılarıyla yapılır; oranı
ağırlıklandırmak gözlemleri bağımsız hale getirmez.
\item Örnek: ``Bu plan, belirtilen dönem ve hedef nüfustaki yetişkinler
için 60 yaş ve üzeri oranını tahmin etmeyi amaçlar. Sonuç, bugünkü
Türkiye nüfusunun oranı veya yaşın herhangi bir sonuç üzerindeki
nedensel etkisi olarak yorumlanamaz.'' Sayısal sonuç ayrıca doğrulanmalıdır.
\end{enumerate}

\textbf{Değerlendirme ölçütleri (10 puan).} Her görev iki ayrı
bileşenden değerlendirilir; bileşen başına doğru ve gerekçeli yanıt
1 puan, eksik veya yanlış yanıt 0 puandır.
\begin{enumerate}
\item Hedef--alt evren--dönem tanımı (1); kapsam değişikliğinin
ağırlık değişikliğinden ayrılması (1).
\item Çerçeve ile yanıtlayıcı dosyasının ayrılması (1); yeniden
seçimin kayıp kapsamı geri getirmediğinin açıklanması (1).
\item İki paydanın doğru yorumu (1); temel ve nihai ağırlık ayrımı (1).
\item Alt evren için tasarım bilgisinin korunması (1); kümeli ve
tabakalı varyans gereğinin açıklanması (1).
\item Kapsamı doğru belirten rapor cümlesi (1); desteklenmeyen
genelleme ve nedensellik iddialarının dışlanması (1).
\end{enumerate}
Doğru formülü yazıp yanlış evrene genellemek tam puan getirmez.

\newcommand{\VtwoN}{9254}
\newcommand{\VtwoAdults}{5569}
\newcommand{\VtwoOlder}{2150}
\newcommand{\VtwoDf}{15}
\newcommand{\VtwoWeighted}{29{,}1516}
\newcommand{\VtwoUnweighted}{38{,}6066}
\newcommand{\VtwoSE}{1{,}5529}
\newcommand{\VtwoLower}{25{,}8416}
\newcommand{\VtwoUpper}{32{,}4615}
\subsection{Sayısal uygulama: dosya oranı ve hedef nüfus oranı}
\label{sec:v2-nhanes-sonuclar}

Kaynakta \VtwoN{} kişi kaydı vardır; bunların \VtwoAdults{} tanesi
20 yaş ve üzerinde, \VtwoOlder{} tanesi 60 yaş ve üzerindedir.
Seçili alanlarda eksiksizlik, pozitif görüşme ağırlıkları ve 15
tabaka--30 PSU yapısı kontrol edilmiştir. Diğer sütunların eksikleri
nedeniyle kayıt silinmemiştir. Dosyanın sayısal kabulü, kamu dosyasının
gizli hane seçim çerçevesini içerdiği anlamına gelmez.

\begin{center}
\begin{tabular}{@{}lr@{}}
\toprule
Özet & Sonuç \\
\midrule
Ağırlıksız yetişkin kayıt oranı & \%\VtwoUnweighted \\
Ağırlıklı alt evren oranı & \%\VtwoWeighted \\
Taylor standart hatası & \VtwoSE{} yüzde puan \\
Tasarım serbestlik derecesi & \VtwoDf \\
Yaklaşık \%95 güven aralığı & [\%\VtwoLower; \%\VtwoUpper] \\
\bottomrule
\end{tabular}
\end{center}

Ağırlıklı pay yaklaşık 69,596 milyon, payda 238,738 milyondur;
bunlar dosyadaki kişi sayısı değil, görüşme ağırlıklarıyla elde edilen
nüfus toplamı tahminleridir. Ağırlıksız oran yanıtlayıcıların bileşimini,
ağırlıklı oran hedef alt nüfusu özetler. Farklı olmaları hesap çelişkisi
veya iki bağımsız araştırma kanıtı değildir.

\textbf{İzlenebilir varyans hesabı.} $B=\sum_iw_id_i$ ve
$z_i=w_id_i(y_i-\widehat p_w)/B$ olsun. Tabaka $h$ içindeki PSU $j$
için $Z_{hj}$ kişi katkılarının toplamıdır. Nihai PSU yaklaşımıyla,
sonlu evren düzeltmesi kullanmadan,
\[
\widehat{\Var}(\widehat p_w)=
\sum_h\frac{m_h}{m_h-1}\sum_j(Z_{hj}-\overline Z_h)^2
\]
hesaplanır. Bütün kayıtlar tasarımda kalır; çocukların alt evren
katkıları sıfırdır. Bu veride yetişkinler bütün 15 tabaka ve 30 PSU'da
bulunduğundan serbestlik derecesi $30-15=15$ alınır
\citep{nchsvariance}. PSU kodları tabaka koduyla birlikte kullanılır.

Aralık $\widehat p_w\pm t_{0{,}975;15}\widehat{\SE}$ biçiminde,
yaklaşık t-Wald öğretim aralığıdır. NCHS'nin bütün oran yayımlama
ölçütlerinin uygulandığı iddia edilmez. Kapsam veya yanıtsızlık
yanlılığı yalnız standart hata hesabıyla ölçülmüş olmaz.

\begin{researchbox}
\textbf{Örnek rapor.} NHANES 2017--2018 görüşme ağırlıkları ve
maskelenmiş varyans tasarımıyla, hedefteki kurum dışı sivil yetişkin
nüfusta 60 yaş ve üzeri oranı yaklaşık \%29,15 tahmin edilmiştir
(Taylor standart hatası 1,55 yüzde puan; yaklaşık \%95 t aralığı
\%25,84--\%32,46). Bu tahmin bugünkü nüfusa, Türkiye'ye veya kurum
nüfusuna doğrudan taşınamaz; nedensel etkiyi ölçmez. CDC'nin belirli
bir resmî tahmininin tekrarı olarak sunulmaz.
\end{researchbox}

\begin{softwarebox}
Tam duyarlıklı sonuçlar \path{companion/vakalar/v2-nhanes/results.json},
30 PSU'nun katkıları aynı klasörde \texttt{psu.csv} içindedir.
\texttt{export.py} çıktıları üretir; \texttt{verify\_delivery.py}
dosya bütünlüğünü, sayısal eşleşmeyi ve kitap bağlantısını denetler.
İkinci hesap pay/payda toplamlarının kovaryansından delta yöntemiyle
aynı varyansa ulaşmıştır. XPT okuyucusu ve t kritik değeri ortaktır;
R/SPSS veya bağımsız bir anket yazılımı çalıştırılmış sayılmaz.
Sayısal eşleşme, örnekleme kapsamının veya yanıt sürecinin kusursuzluğunu
kanıtlamaz.
\end{softwarebox}